\documentclass[11pt]{amsart}
\usepackage[margin=1.1in]{geometry}
\usepackage{amssymb,booktabs,hyperref}
\hypersetup{colorlinks=true,linkcolor=blue,citecolor=blue,urlcolor=blue}

\newtheorem{theorem}{Theorem}[section]
\newtheorem{proposition}[theorem]{Proposition}
\newtheorem{lemma}[theorem]{Lemma}
\newtheorem{corollary}[theorem]{Corollary}
\theoremstyle{definition}
\newtheorem{definition}[theorem]{Definition}
\newtheorem{remark}[theorem]{Remark}
\numberwithin{equation}{section}

\newcommand{\R}{\mathbb{R}}
\newcommand{\C}{\mathbb{C}}
\newcommand{\Herm}{\mathrm{Herm}_n}
\newcommand{\Pn}{\mathcal{P}_n}
\newcommand{\Dn}{\mathcal{D}_n}
\newcommand{\Tr}{\operatorname{Tr}}
\newcommand{\tr}{\operatorname{tr}}
\newcommand{\Ric}{\operatorname{Ric}}
\newcommand{\hs}{\mathrm{hs}}
\newcommand{\ph}{\hat\varphi}
\newcommand{\eps}{\varepsilon}

\title[The WYD curvature conjecture]{Scalar curvature of the Wigner--Yanase--Dyson metrics:\\ closed form and monotonicity under mixing}
\author{Adam Bornemann}
\author{Claude (Anthropic)}
\date{25 November 2026. Working document; not independently reviewed.}

\begin{document}

\begin{abstract}
For the Wigner--Yanase--Dyson metric $g^\beta$, $\beta\in\R\setminus\{0,1\}$, on the faithful $n\times n$ density
matrices we prove a closed form for the scalar curvature as a sum over triples of eigenvalues of one explicit symmetric
function $\hs_\beta$ (Theorem~\ref{thm:closed}). We then prove that strict monotonicity under mixing
(Gibilisco--Isola, Conjecture~8.1, at the parameter $\beta$) follows from Dittmann's inequalities (62)--(65) for
$\hs_\beta$, with (62) strict (Theorem~\ref{thm:reduction}). We prove (62) for every $\beta\ne\frac12$ (Theorem~\ref{thm:62}) and (63) for $\beta\in[-1,2]$
(Theorem~\ref{thm:63}) by hand, and (64) and (65) for $\beta\in[-1,2]$ (Theorem~\ref{thm:6465}) by a finite certificate in integer arithmetic: the numerators are exponential polynomials in two gaps, and their
Newton coefficients on their own exponents are polynomials in one parameter with nonnegative integer coefficients. So the conjecture holds for every $\beta\in[-1,2]\setminus\{0,\frac12,1\}$ and every $n$
(Corollary~\ref{cor:main}). The whole proof is machine-checked in Lean (Section~\ref{sec:status}).
\end{abstract}

\maketitle

\section{Setting and statements}

\subsection*{The metric} $\Herm$ denotes the real vector space of Hermitian $n\times n$ matrices ($n\ge 2$), $\Pn\subset\Herm$ the
open cone of positive definite ones and $\Dn=\{\rho\in\Pn:\Tr\rho=1\}$. Fix $\beta\in\R\setminus\{0,1\}$. For $\rho=U\operatorname{diag}(\lambda)U^*$
and $X\in\Herm$ put $\tilde X=U^*XU$ and
\begin{equation}\label{eq:metric}
g^\beta_\rho(X,X)=\sum_{i,j}c_\beta(\lambda_i,\lambda_j)\,|\tilde X_{ij}|^2,\qquad
c_\beta(x,y)=\frac{(x^\beta-y^\beta)(x^{1-\beta}-y^{1-\beta})}{\beta(1-\beta)(x-y)^2},\quad c_\beta(x,x)=\frac1x .
\end{equation}
Since $(x^\beta-y^\beta)/\beta$ and $(x^{1-\beta}-y^{1-\beta})/(1-\beta)$ have the sign of $x-y$, $c_\beta>0$, so $g^\beta$ is a Riemannian
metric (real-analytic, by Lemma~\ref{lem:hessian}); it is invariant under $\rho\mapsto U\rho U^*$ and under $\beta\leftrightarrow1-\beta$.
It is a monotone metric for $-1\le\beta\le2$ (Hasegawa--Petz \cite{hp}), but monotonicity of the metric is not used
below; nor is the conjecture's restriction to $p=1/\beta$ near $1$ \cite{gi}. The limits $\beta\to0,1$ give the Kubo--Mori metric, for which the statement is a theorem (Petz's conjecture, proved in
Lean in this repository); we exclude them.

\subsection*{Conventions} $R(X,Y)Z=\nabla_X\nabla_YZ-\nabla_Y\nabla_XZ-\nabla_{[X,Y]}Z$, $\Ric(Y,Z)=\tr(X\mapsto R(X,Y)Z)$, and $S$ is the
trace of $\Ric$; the unit sphere of dimension $N$ has $S=N(N-1)$. $S_\beta(\rho)$ is the scalar curvature of $(\Dn,g^\beta)$ at $\rho$. By unitary
invariance it depends only on the eigenvalues, $S_\beta(\rho)=S_\beta(\lambda)$.

For $x,y\in\R^n$, $x$ is \emph{majorized by} $y$, $x\prec y$, if $\sum x_i=\sum y_i$ and, with the entries of both sorted increasingly,
$\sum_{i\le k}x_i\ge\sum_{i\le k}y_i$ for every $k$. A state $\rho$ is \emph{more mixed} than $\varrho$ if its eigenvalue vector is majorized by that
of $\varrho$.

\subsection*{Notation for divided differences} For a smooth $f$ on $(0,\infty)$, $f^{[1]}(x,y)$ and $f^{[2]}(x,y,z)$ are its first and second divided
differences; by the Hermite--Genocchi formula
\begin{equation}\label{eq:HG}
f^{[k]}(x_0,\dots,x_k)=\int_{\Delta_k}f^{(k)}(t_0x_0+\dots+t_kx_k)\,dt,
\end{equation}
so they are smooth symmetric functions of positive arguments, coincident ones included, with $f^{[1]}(x,x)=f'(x)$ and $f^{[2]}(x,x,x)=\frac12f''(x)$.

\section{The WYD metric is a Hessian metric}

Put
\begin{equation}
p=\frac1\beta,\qquad \varphi(\mu)=\frac{\mu^{p-1}}{\beta(1-\beta)},\qquad \Phi(u)=\frac{\Tr u^{p}}{1-\beta}\quad(u\in\Pn).
\end{equation}
Then $\varphi(\lambda^\beta)=\lambda^{1-\beta}/(\beta(1-\beta))$, and, since $p-1=(1-\beta)/\beta$,
\begin{equation}\label{eq:phiprime}
\varphi'(\mu)=\frac{\mu^{p-2}}{\beta^2}>0 ,\qquad \varphi''(\mu)=\frac{1-2\beta}{\beta\mu}\,\varphi'(\mu).
\end{equation}

\begin{lemma}\label{lem:dk}
Let $f$ be holomorphic on $\C\setminus(-\infty,0]$ and $u=\operatorname{diag}(\mu)\in\Pn$. Then
$\frac{d}{dt}\big|_0\Tr F(u+tA)=\Tr(f(u)A)$ if $F'=f$, and
\[
(Df(u)[A])_{ij}=f^{[1]}(\mu_i,\mu_j)A_{ij},\qquad
(D^2f(u)[A,B])_{ik}=\sum_jf^{[2]}(\mu_i,\mu_j,\mu_k)(A_{ij}B_{jk}+B_{ij}A_{jk}).
\]
\end{lemma}
\begin{proof}
For $\Gamma$ a contour in $\{\operatorname{Re}z>0\}$ around the $\mu_i$ and $R=(z-u)^{-1}$, $f(u)=\frac1{2\pi i}\oint_\Gamma f(z)R\,dz$. As $DR[A]=RAR$, we get
\[Df(u)[A]=\frac1{2\pi i}\oint_\Gamma f\,RAR\,dz,\qquad D^2f(u)[A,B]=\frac1{2\pi i}\oint_\Gamma f\,(RARBR+RBRAR)\,dz .\]
In coordinates the scalar integrals are
$\frac1{2\pi i}\oint f(z)\prod_{l=0}^{k}(z-x_l)^{-1}dz=f^{[k]}(x_0,\dots,x_k)$ (Cauchy). For the trace, $\Tr(F^{[1]}(\mu_i,\mu_j)A_{ij})_{ij}=\sum_if(\mu_i)A_{ii}$.
\end{proof}

\begin{lemma}\label{lem:hessian}
The map $\rho\mapsto u=\rho^\beta$ is a real-analytic diffeomorphism of $\Pn$ onto itself, and it carries $g^\beta$ (extended to $\Pn$ by
\eqref{eq:metric}) to the Hessian metric $G_u=D^2\Phi(u)$. It carries $\Dn$ to the level set $\{\Phi=1/(1-\beta)\}$.
\end{lemma}
\begin{proof}
The inverse is $u\mapsto u^{1/\beta}$, and both are analytic. By Lemma~\ref{lem:dk},
\[D\Phi(u)[A]=\frac{p}{1-\beta}\Tr(u^{p-1}A)=\Tr(\varphi(u)A),\qquad D^2\Phi(u)[A,B]=\Tr(D\varphi(u)[A]\,B).\] With $A=B=D\rho^\beta[X]$ and $\varphi(\rho^\beta)=\rho^{1-\beta}/(\beta(1-\beta))$, the chain rule gives
\[D^2\Phi[A,A]=\frac{1}{\beta(1-\beta)}\Tr(D\rho^{1-\beta}[X]\,D\rho^\beta[X]).\]
At $\rho=\operatorname{diag}(\lambda)$, Lemma~\ref{lem:dk} turns this into
$\sum_{ij}c_\beta(\lambda_i,\lambda_j)|X_{ij}|^2$ (for $\lambda_i=\lambda_j$ the factor is $\lambda_i^{-\beta}\lambda_i^{\beta-1}=1/\lambda_i$). The case of general $\rho$ follows by unitary
invariance of both sides. Finally $\Tr\rho=\Tr u^p=(1-\beta)\Phi(u)$.
\end{proof}

\noindent The next lemma is standard (Shima \cite{shima}); Lemma~\ref{lem:dk} is the Daleckii--Krein formula \cite{bhatia}. Both are proved to fix signs and normalisations.

\begin{lemma}\label{lem:hesscurv}
Let $G=D^2\Phi$ be a Riemannian Hessian metric on an open subset of a finite-dimensional vector space, $T=D^3\Phi$, and $\tau(Z)=\sum_aT(e_a,e_a,Z)$ for a $G$-orthonormal
basis $(e_a)$. Let $T_XY$ be defined by $G(T_XY,Z)=T(X,Y,Z)$. Then, for constant vector fields,
\begin{equation*}
\nabla_XY=\tfrac12T_XY,\qquad R(X,Y)=-\tfrac14[T_X,T_Y],\qquad S=\tfrac14\Big(\sum_{a,b,c}T(e_a,e_b,e_c)^2-\sum_a\tau(e_a)^2\Big).
\end{equation*}
\end{lemma}
\begin{proof}
$DG=T$ is totally symmetric, so Koszul's formula gives $2G(\nabla_XY,Z)=T(X,Y,Z)$. Differentiating $G(T_YZ,W)=T(Y,Z,W)$ gives
\[G(D_X(T_YZ),W)=D^4\Phi(X,Y,Z,W)-T(X,T_YZ,W).\]
Hence
\[R(X,Y)Z=\tfrac12\big(D_X(T_YZ)-D_Y(T_XZ)\big)+\tfrac14(T_XT_YZ-T_YT_XZ)=\big(-\tfrac12+\tfrac14\big)[T_X,T_Y]Z ,\]
the fourth derivatives cancelling. Then
\[\Ric(Y,Z)=-\frac14\sum_aG([T_{e_a},T_Y]Z,e_a)=\frac14\Big(\sum_aG(T_Ye_a,T_Ze_a)-\tau(T_YZ)\Big),\]
using the symmetry of $T$:
\[G(T_{e_a}T_YZ,e_a)=G(T_YZ,T_{e_a}e_a),\qquad G(T_YT_{e_a}Z,e_a)=G(T_Ze_a,T_Ye_a).\]
 Tracing over $Y=Z=e_b$ and using $\sum_bT_{e_b}e_b=\sum_a\tau(e_a)e_a$ gives $S$.
\end{proof}

\begin{proposition}\label{prop:gauss}
Let $S^{\mathrm{cone}}$ be the scalar curvature of $(\Pn,G)$. At every point of $\Dn$,
\[S_\beta=S^{\mathrm{cone}}+\frac{(n^2-1)(n^2-2)}{4}.\]
\end{proposition}
\begin{proof}
$\Phi$ is homogeneous of degree $p$. Let $E_u=u$ be the Euler field. Euler's relation for $D\Phi$ and $D^2\Phi$ gives
\[G(E,V)=(p-1)D\Phi(u)[V],\qquad T(E,V,W)=(p-2)G(V,W).\]
Hence $E$ is normal to the level sets of $\Phi$, and $G(E,E)=p(p-1)\Phi$. On $\Dn$, $\Phi=1/(1-\beta)=p/(p-1)$, so $G(E,E)=p^2$ and $\nu=E/p$ is a unit normal
along $\Dn$ (its sign is irrelevant). For every $V$,
\[\nabla_VE=D_VE+\tfrac12T_VE=V+\tfrac12(p-2)V=\tfrac p2V .\]
Two consequences. First, $\nabla_V\nu=\frac12V$ for $V$ tangent to $\Dn$: the hypersurface is totally umbilic with all principal curvatures equal to
$\frac12$. Second, for any vector fields, $R(X,Y)E=\frac p2(\nabla_XY-\nabla_YX-[X,Y])=0$, so $\Ric(E,E)=0$. The Gauss equation, traced twice, reads
$S_\beta=S^{\mathrm{cone}}-2\Ric(\nu,\nu)+(\tr A)^2-|A|^2$ with $A=\frac12\mathrm{Id}$ on a space of dimension $N=n^2-1$, and $\frac{N^2}4-\frac N4=\frac{(n^2-1)(n^2-2)}4$.
\end{proof}

\section{The closed form}

For $x,y,z>0$ put
\begin{align}
w(x,y)&=\varphi^{[1]}(x^\beta,y^\beta)>0,\qquad F(x,y,z)=\varphi^{[2]}(x^\beta,y^\beta,z^\beta),\\
h(x,y,z)&=\frac{F(x,y,z)^2}{2\,w(x,y)w(y,z)w(z,x)}-\frac{F(x,x,y)F(x,x,z)}{w(x,x)w(x,y)w(x,z)},\label{eq:h}\\
\hs_\beta(x,y,z)&=\tfrac13\big(h(x,y,z)+h(y,z,x)+h(z,x,y)\big).\label{eq:hs}
\end{align}
By \eqref{eq:HG} and \eqref{eq:phiprime} these are smooth on $(0,\infty)^3$, and $\hs_\beta$ is symmetric ($h$ is symmetric in its last two arguments). Scaling $\lambda\mapsto c\lambda$
multiplies $w$ by $c^{1-2\beta}$ and $F$ by $c^{1-3\beta}$, so $\hs_\beta$ is homogeneous of degree $2(1-3\beta)-3(1-2\beta)=-1$.

\begin{theorem}\label{thm:closed}
For every $\beta\in\R\setminus\{0,1\}$ and every $\rho\in\Dn$ with eigenvalues $\lambda_1,\dots,\lambda_n$,
\begin{equation}\label{eq:closed}
S_\beta(\rho)=\sum_{i,j,k}\hs_\beta(\lambda_i,\lambda_j,\lambda_k)-\sum_i\hs_\beta(\lambda_i,\lambda_i,\lambda_i)+\frac{(n^2-1)(n^2-2)}4 .
\end{equation}
\end{theorem}
\begin{proof}
By unitary invariance take $\rho$ diagonal; $u=\operatorname{diag}(\mu)$, $\mu_i=\lambda_i^\beta$. Write $w_{ij}$, $F_{ijk}$ for $w,F$ at the corresponding eigenvalues. By Lemma~\ref{lem:dk},
$G(A,B)=\sum_{ij}w_{ij}A_{ij}B_{ji}$ and
\[T(A,B,C)=C_3(A,B,C)+C_3(B,A,C),\qquad C_3(A,B,C)=\sum_{i,j,k}F_{ijk}A_{ij}B_{jk}C_{ki}.\]
\emph{Completeness.} Let $(e_a)$ be a $G$-orthonormal basis of $\Herm$. It is also a $\C$-basis of $M_n(\C)$, orthonormal for the complex bilinear
extension of $G$, so $A=\sum_aG(A,e_a)e_a$ for all $A\in M_n(\C)$. For $A=E_{rs}$ this says
\begin{equation}\label{eq:complete}
\sum_a(e_a)_{pq}(e_a)_{sr}=\frac{\delta_{pr}\delta_{qs}}{w_{rs}} .
\end{equation}
\emph{$|T|^2$.} Total symmetry of $T$ and relabelling give $\sum_{abc}T(e_a,e_b,e_c)^2=2\Sigma_1+2\Sigma_2$ with $\Sigma_1=\sum C_3(e_a,e_b,e_c)^2$ and
$\Sigma_2=\sum C_3(e_a,e_b,e_c)C_3(e_b,e_a,e_c)$ (the substitution $a\leftrightarrow b$ exchanges the two squares). In $\Sigma_1$, with primed indices for the second
factor, \eqref{eq:complete} forces $i=j'$, $j=i'$; $j=k'$, $k=j'$; $k=i'$, $i=k'$: all six indices coincide, and $\Sigma_1=\sum_iF_{iii}^2/w_{ii}^3$. In $\Sigma_2$ the second factor is
$\sum F_{i'j'k'}(e_b)_{i'j'}(e_a)_{j'k'}(e_c)_{k'i'}$, and \eqref{eq:complete} forces $i=k'$, $j=j'$, $k=i'$, so by symmetry of $F$,
$\Sigma_2=\sum_{ijk}F_{ijk}^2/(w_{ij}w_{jk}w_{ki})$.

\emph{$|\tau|^2$.} $T(e_a,e_a,Z)=2C_3(e_a,e_a,Z)$, and $\sum_a(e_a)_{ij}(e_a)_{jk}=\delta_{ik}/w_{ij}$, so $\tau(Z)=2\sum_it_iZ_{ii}$ with $t_i=\sum_jF_{iij}/w_{ij}$; then
$\sum_a\tau(e_a)^2=4\sum_it_i^2/w_{ii}$.

By Lemma~\ref{lem:hesscurv},
\[S^{\mathrm{cone}}=\sum_{ijk}\frac{F_{ijk}^2}{2w_{ij}w_{jk}w_{ki}}-\sum_{ijk}\frac{F_{iij}F_{iik}}{w_{ii}w_{ij}w_{ik}}+\sum_i\frac{F_{iii}^2}{2w_{ii}^3}
=\sum_{ijk}h(\lambda_i,\lambda_j,\lambda_k)-\sum_ih(\lambda_i,\lambda_i,\lambda_i),\]
because $h(x,x,x)=-F(x,x,x)^2/(2w(x,x)^3)$. A sum over all triples does not change under cyclic relabelling, so $h$ may be replaced by $\hs_\beta$.
Proposition~\ref{prop:gauss} finishes the proof.
\end{proof}

\section{Reduction to Dittmann's inequalities}

Write $\hs'$ for the partial derivative of $\hs=\hs_\beta$ in its first argument. Dittmann's inequalities are, for all $0<x<y$ and $l,m>0$,
\begin{align}
0&\le 2\hs'(x,x,y)-\hs'(y,x,x)-2\hs'(y,x,y)+\hs'(x,y,y),\tag{62}\label{62}\\
0&\le \hs'(x,x,l)-\hs'(y,y,l),\tag{63}\label{63}\\
0&\le \hs'(x,y,l)-\hs'(y,x,l),\tag{64}\label{64}\\
0&\le \hs'(x,l,m)-\hs'(y,l,m).\tag{65}\label{65}
\end{align}
Dittmann's proof for the Kubo--Mori metric uses nothing about his function beyond smoothness and symmetry, and the same is true here.

\begin{theorem}\label{thm:reduction}
Let $\hs$ be any symmetric $C^1$ function on $(0,\infty)^3$ satisfying \eqref{62}--\eqref{65}, and
$\mathcal S(\lambda)=\sum_{ijk}\hs(\lambda_i,\lambda_j,\lambda_k)-\sum_i\hs(\lambda_i,\lambda_i,\lambda_i)$ for $\lambda\in(0,\infty)^n$. If $x\prec y$ then $\mathcal S(y)\le\mathcal S(x)$. If moreover \eqref{62} is
strict and the sorted vectors differ, then $\mathcal S(y)<\mathcal S(x)$.

Consequently, if $\hs_\beta$ satisfies \eqref{62}--\eqref{65} with \eqref{62} strict, then $S_\beta(\varrho)<S_\beta(\rho)$ whenever $\rho\in\Dn$ is more mixed than $\varrho\in\Dn$ and
their spectra differ.
\end{theorem}
\begin{proof}
\emph{Gradient.} By symmetry of $\hs$, $\partial_a\mathcal S=3\big(\sum_{jk}\hs'(\lambda_a,\lambda_j,\lambda_k)-\hs'(\lambda_a,\lambda_a,\lambda_a)\big)$. Fix $a\neq b$ with
$\lambda_a<\lambda_b$. Split the double sums according to whether $j$ and $k$ lie in $\{a,b\}$. Writing $\hs'_{a|jk}$ for $\hs'(\lambda_a,\lambda_j,\lambda_k)$, the terms with both
in $\{a,b\}$, together with $-\hs'_{a|aa}+\hs'_{b|bb}$, sum to
$2\hs'_{a|ab}+\hs'_{a|bb}-\hs'_{b|aa}-2\hs'_{b|ab}$. So
\begin{multline}\label{eq:dittmann}
\tfrac13(\partial_a-\partial_b)\mathcal S=\big(2\hs'_{a|ab}-\hs'_{b|aa}-2\hs'_{b|ab}+\hs'_{a|bb}\big)\\
+2\sum_{r\ne a,b}\big((\hs'_{a|ar}-\hs'_{b|br})+(\hs'_{a|br}-\hs'_{b|ar})\big)+\sum_{j,k\ne a,b}(\hs'_{a|jk}-\hs'_{b|jk}).
\end{multline}
The groups are $\ge0$ by \eqref{62}, by \eqref{63} and \eqref{64}, and by \eqref{65}. Hence $\partial_a\mathcal S\ge\partial_b\mathcal S$ when $\lambda_a<\lambda_b$, strictly if \eqref{62} is
strict; and $\partial_a\mathcal S=\partial_b\mathcal S$ when $\lambda_a=\lambda_b$.

\emph{Schur--Ostrowski.} $\mathcal S$ is invariant under permutations, so let $x,y$ be sorted increasingly. Put $v=x-y$ and $V_i=\sum_{j\le i}v_j$; then $V_i\ge0$ and $V_n=0$. The points
$z(t)=y+tv=(1-t)y+tx$, $t\in[0,1]$, are positive and sorted, so $\partial_i\mathcal S\ge\partial_{i+1}\mathcal S$ there and, by Abel summation,
\[\frac{d}{dt}\mathcal S(z(t))=\sum_i\partial_i\mathcal S\,v_i=\sum_{i<n}(\partial_i\mathcal S-\partial_{i+1}\mathcal S)V_i\ge0 .\]
For the strict statement let $x\ne y$ and $t<1$. Abel summation also gives $x\cdot v=-\sum_{i<n}(x_{i+1}-x_i)V_i\le0$, hence
\[-\sum_{i<n}(z_{i+1}-z_i)V_i=z(t)\cdot v=x\cdot v-(1-t)|v|^2<0 .\]
So for some $i$ both $V_i>0$ and $z_i<z_{i+1}$; the $i$-th term of the derivative is then strictly positive. Thus $\mathcal S(z(t))$ is strictly increasing on $[0,1)$,
and $\mathcal S(y)<\mathcal S(x)$.

The last claim is Theorem~\ref{thm:closed}.
\end{proof}

\begin{remark}
The reduction is machine-checked for the Kubo--Mori function in the Lean file
\texttt{Dittmann/Mixing.lean}. The Schur--Ostrowski criterion is machine-checked for general functions, in its weak and strict forms:
\texttt{le\_of\_isMajorizedBy} and \texttt{lt\_of\_isMajorizedBy} in
\texttt{ForMathlib/Majorization.lean}. The proof above is the same, following \cite{dittmann}.
\end{remark}

\section{Explicit forms of \texorpdfstring{$\hs_\beta$}{hs}}\label{sec:forms}

Put $P_i=\lambda_i^\beta$, $Q_i=\lambda_i^{1-\beta}$ (so $\lambda_i=P_iQ_i$), $D_{ij}=(P_i-P_j)(Q_i-Q_j)$, $A=\det\begin{pmatrix}1&1&1\\P_1&P_2&P_3\\Q_1&Q_2&Q_3\end{pmatrix}$ and
\[\kappa_{ij}=\frac{(1-\beta)Q_i}{Q_j-Q_i}-\frac{\beta P_i}{P_j-P_i}\qquad(\lambda_i\ne\lambda_j).\]

\begin{proposition}\label{prop:kappa}
For pairwise distinct arguments
\begin{equation}\label{eq:kappaform}
\hs_\beta(\lambda_1,\lambda_2,\lambda_3)=\frac{\beta(1-\beta)A^2}{2D_{12}D_{23}D_{31}}-\frac13\sum_{a}\frac{\kappa_{ab}\kappa_{ac}}{\lambda_a}\qquad(\{a,b,c\}=\{1,2,3\}),
\end{equation}
and for $\lambda_a\neq\lambda_b$
\begin{equation}\label{eq:twopoint}
\hs_\beta(\lambda_a,\lambda_a,\lambda_b)=\frac1{\lambda_a}\Big(\frac{\kappa_{ab}^2}2+\frac{1-2\beta}3\kappa_{ab}\Big)-\frac{\kappa_{ba}^2}{3\lambda_b}.
\end{equation}
\end{proposition}
\begin{proof}
From $\varphi(P)=Q/(\beta(1-\beta))$ on the curve $(P_i,Q_i)$ and \eqref{eq:phiprime},
\[w_{ij}=\frac{1}{\beta(1-\beta)}\,\frac{Q_i-Q_j}{P_i-P_j},\qquad w_{ii}=\frac{Q_i}{\beta^2P_i},\qquad F_{iij}=\frac{w_{ii}-w_{ij}}{P_i-P_j}.\]
Hence
\begin{equation}\label{eq:Fw}
\frac{F_{iij}}{w_{ij}}=\frac1{P_i-P_j}\Big(\frac{w_{ii}}{w_{ij}}-1\Big)=\frac{1-\beta}{\beta P_i}\,\frac{Q_i}{Q_i-Q_j}-\frac1{P_i-P_j}=-\frac{\kappa_{ij}}{\beta P_i},
\end{equation}
and therefore $\dfrac{F_{aab}F_{aac}}{w_{aa}w_{ab}w_{ac}}=\dfrac{\kappa_{ab}\kappa_{ac}}{\beta^2P_a^2w_{aa}}=\dfrac{\kappa_{ab}\kappa_{ac}}{\lambda_a}$, since $\beta^2P_a^2w_{aa}=P_aQ_a$. The second divided difference is
$F_{123}=\dfrac{A}{\beta(1-\beta)\prod_{i<j}(P_j-P_i)}$, and $w_{12}w_{23}w_{31}=(\beta(1-\beta))^{-3}\prod_{i<j}\frac{Q_i-Q_j}{P_i-P_j}$, so
$\dfrac{F_{123}^2}{2w_{12}w_{23}w_{31}}=\dfrac{\beta(1-\beta)A^2}{2D_{12}D_{23}D_{31}}$ (the product $D_{12}D_{23}D_{31}$ equals $\prod_{i<j}(P_i-P_j)(Q_i-Q_j)$). This is \eqref{eq:kappaform}.

For \eqref{eq:twopoint}, $3\hs(\lambda_a,\lambda_a,\lambda_b)=2h(\lambda_a,\lambda_a,\lambda_b)+h(\lambda_b,\lambda_a,\lambda_a)$. By \eqref{eq:Fw} the first term of $h$ is, for every ordering,
$\dfrac{F_{aab}^2}{2w_{aa}w_{ab}^2}=\dfrac{\kappa_{ab}^2}{2\lambda_a}$. The second term of $h(\lambda_b,\lambda_a,\lambda_a)$ is $\kappa_{ba}^2/\lambda_b$. In $h(\lambda_a,\lambda_a,\lambda_b)$ it is
$\dfrac{F_{aaa}}{w_{aa}^2}\cdot\dfrac{F_{aab}}{w_{ab}}$ with $F_{aaa}=\frac12\varphi''(P_a)=\frac{1-2\beta}{2\beta P_a}w_{aa}$, i.e.\
$-\dfrac{(1-2\beta)\kappa_{ab}}{2\beta^2P_a^2w_{aa}}=-\dfrac{(1-2\beta)\kappa_{ab}}{2\lambda_a}$. Adding up gives \eqref{eq:twopoint}.
\end{proof}

\subsection*{The hyperbolic form} From now on
\[d=2\beta-1,\qquad D=d^2,\qquad \lambda_i=e^{2\sigma_i},\qquad X_i=e^{-2\sigma_i}=1/\lambda_i,\]
so $\beta=\frac{1+d}2$, $P_i=e^{(1+d)\sigma_i}$, $Q_i=e^{(1-d)\sigma_i}$. For $d\notin\{0,\pm1\}$ let
\begin{equation}
\ph(s)=\frac{\sinh s-\frac1d\sinh ds}{1-d^2},\qquad K=\ph/\ph',\qquad \ph'(s)=\frac{\cosh s-\cosh ds}{1-d^2}.
\end{equation}
$\ph$ is odd and $\ph''=d^2\ph+\sinh s$. Its Taylor series is $\sum_{j\ge1}(1+D+\dots+D^{j-1})s^{2j+1}/(2j+1)!$, which is entire in $(s,D)$ and extends the definition to $d\in\{0,\pm1\}$. For $D\ge0$
all the coefficients are $\ge0$ and the first is positive, so $\ph,\ph',\ph''>0$ for $s>0$. $K$ is odd; near $0$ it is analytic, $K(s)=s/3+\dots$.

\begin{proposition}\label{prop:H}
For pairwise distinct $\sigma_i$ and $d\notin\{0,\pm1\}$, $\hs_\beta=d^2H(\sigma;D)$ with
\begin{equation}\label{eq:H}
H=\frac{\big(\sum_{\mathrm{cyc}}e^{-\sigma_k}\ph(\sigma_i-\sigma_j)\big)^2}{16\,\ph'(\sigma_1-\sigma_2)\ph'(\sigma_2-\sigma_3)\ph'(\sigma_3-\sigma_1)}
-\frac1{12}\sum_aX_a\big(K(\sigma_b-\sigma_a)-1\big)\big(K(\sigma_c-\sigma_a)-1\big),
\end{equation}
the first sum over the cyclic $(i,j,k)$. Moreover
\begin{equation}\label{eq:kappaK}
\kappa_{ab}=\tfrac d2\big(1-K(\sigma_b-\sigma_a)\big),\qquad \kappa_{ab}+\kappa_{ba}=d .
\end{equation}
For $d\ne0,\pm1$, $H$ extends smoothly to coincident $\sigma_i$, since $\hs_\beta$ is smooth.
\end{proposition}
\begin{proof}
Let $\delta=\sigma_i-\sigma_j$. Then $D_{ij}=4e^{\sigma_i+\sigma_j}\sinh\frac{(1+d)\delta}2\sinh\frac{(1-d)\delta}2=2(1-d^2)e^{\sigma_i+\sigma_j}\ph'(\delta)$. Next,
$A=\sum_{\mathrm{cyc}}(P_iQ_j-P_jQ_i)=2\sum_{\mathrm{cyc}}e^{\sigma_i+\sigma_j}\sinh(d\delta_{ij})$. Insert $\sinh(d\delta)=d\sinh\delta-d(1-d^2)\ph(\delta)$ and use
$\sum_{\mathrm{cyc}}e^{\sigma_i+\sigma_j}\sinh\delta_{ij}=\frac12\sum_{\mathrm{cyc}}(e^{2\sigma_i}-e^{2\sigma_j})=0$ to get $A=-2d(1-d^2)\sum_{\mathrm{cyc}}e^{\sigma_i+\sigma_j}\ph(\delta_{ij})$. With
$\beta(1-\beta)=\frac{1-d^2}4$ the first term of \eqref{eq:kappaform} becomes
\[\frac{(1-d^2)}{8}\cdot\frac{4d^2(1-d^2)^2\big(\sum e^{\sigma_i+\sigma_j}\ph\big)^2}{8(1-d^2)^3e^{2(\sigma_1+\sigma_2+\sigma_3)}\prod\ph'}=\frac{d^2}{16}\,\frac{\big(\sum e^{-\sigma_k}\ph(\delta_{ij})\big)^2}{\prod\ph'} .\]
For \eqref{eq:kappaK}, let $s=\sigma_b-\sigma_a$. With $\frac1{e^x-1}=\frac12(\coth\frac x2-1)$,
$\kappa_{ab}=\frac{(1-d)/2}{e^{(1-d)s}-1}-\frac{(1+d)/2}{e^{(1+d)s}-1}=\frac d2+\frac14\big[(1-d)\coth\frac{(1-d)s}2-(1+d)\coth\frac{(1+d)s}2\big]$, and the bracket, over its common denominator
$\sinh\frac{(1+d)s}2\sinh\frac{(1-d)s}2=\frac12(\cosh s-\cosh ds)$, has numerator $\frac12[(1-d)(\sinh s+\sinh ds)-(1+d)(\sinh s-\sinh ds)]=\sinh ds-d\sinh s$. So the bracket is
$-2d\,\ph(s)/\ph'(s)$. The sum rule follows as $K$ is odd. Inserting $\kappa_{ab}\kappa_{ac}=\frac{d^2}4(K-1)(K-1)$ gives the second term.
\end{proof}

\begin{proposition}[translation]\label{prop:translate}
Let $d\ne0,\pm1$. Then:
\begin{enumerate}
\item \eqref{62} for all $x<y$ (strictly) $\iff$ the qubit curvature $S_\beta(\lambda,1-\lambda)$ has $\lambda$-derivative $>0$ for $0<\lambda<\frac12$.
\item \eqref{63} $\iff$ $(e^{2s}\psi'(s))'\le0$ for all $s\in\R$, with $\psi=e^{2s}(K-1)(3K+1)-2(K+1)^2$; i.e.\ $2\psi'+\psi''\le0$.
\item \eqref{64} $\iff$ $(e^{-2\sigma_1}\partial_{\sigma_1}-e^{-2\sigma_2}\partial_{\sigma_2})H\ge0$ for $\sigma_1<\sigma_2$ and all $\sigma_3$.
\item \eqref{65} $\iff$ $(\partial_{\sigma_1}^2-2\partial_{\sigma_1})H\le0$ for all $\sigma$.
\end{enumerate}
In (3) and (4) it is enough to check pairwise distinct $\sigma_i$.
\end{proposition}
\begin{proof}
(1) The right side of \eqref{62} is $\frac13(\partial_x-\partial_y)\mathcal P(x,y)$ with $\mathcal P(x,y)=3\hs(x,x,y)+3\hs(x,y,y)$, since $\partial_x\hs(x,x,y)=2\hs'(x,x,y)$ and so on. By
\eqref{eq:closed}, $S_\beta(\lambda,1-\lambda)=\mathcal P(\lambda,1-\lambda)+\frac32$. $(\partial_x-\partial_y)\mathcal P$ is homogeneous of degree $-2$, so its sign on $\{x<y\}$ is determined on $x+y=1$.

(2) Let $q(x)=\hs(x,x,l)$. Then $q'(x)=2\hs'(x,x,l)$, so \eqref{63} says that $q$ is concave. Put $l=e^{2s}x$. By \eqref{eq:twopoint}, \eqref{eq:kappaK} and $1-2\beta=-d$ (here $a$ is $x$, $b$ is $l$),
\[q=\frac{e^{2s}}l\Big(\frac{d^2}8(1-K)^2-\frac{d^2}6(1-K)\Big)-\frac{d^2}{12\,l}(1+K)^2=\frac{d^2}{24\,l}\psi(s),\]
because $3(1-K)^2-4(1-K)=(K-1)(3K+1)$. Since $x=le^{-2s}$, $\frac d{dx}=-\frac{e^{2s}}{2l}\frac d{ds}$ and $q''(x)=\frac{d^2}{96\,l^3}e^{2s}(e^{2s}\psi')'$.

(3), (4) Use $\partial_\lambda=\frac12e^{-2\sigma}\partial_\sigma$, so that $\partial_\lambda^2=\frac14e^{-4\sigma}(\partial_\sigma^2-2\partial_\sigma)$, and $d^2>0$. \eqref{64} says $(\partial_x-\partial_y)\hs(x,y,l)\ge0$, and
\eqref{65} says that $\hs'(\cdot,l,m)$ is nonincreasing, i.e.\ $\partial_x^2\hs\le0$. Distinct arguments are dense and $\hs$ is smooth.
\end{proof}

\subsection*{The per-gap form used in the tails} In this subsection $y,g,G,Y,p,m$ denote per-gap quantities (not the objects of Sections 2--4 with the same letters). For $k\ne 0$ let $B(k,s)=k/(e^{ks}-1)$ and $B(0,s)=1/s$. For a gap $s>0$ put
\[y=B(1-d,s),\quad g=y+1-d,\quad G=B(1+d,s),\quad Y=G+1+d,\quad p=y-G,\quad m=Y-g .\]

\begin{proposition}\label{prop:gy}
Let $\sigma_a<\sigma_b<\sigma_c$, with gaps $u=\sigma_b-\sigma_a$, $v=\sigma_c-\sigma_b$, $w=u+v$, and $d\ne0,\pm1$. Put $\Pi=Y_u+G_v$, $\Theta=g_u+y_v$. Then $p=d(1-K)$, $m=d(1+K)$ for each gap;
\begin{equation}\label{eq:gy}
96\,d^2H=\frac{12X_b\,(Y_uy_v-g_uG_v)^2}{\Pi\,\Theta}-8\big(X_ap_up_w+X_bm_up_v+X_cm_vm_w\big);
\end{equation}
$y_w=y_uy_v/\Theta$, $g_w=g_ug_v/\Theta$, $Y_w=Y_uY_v/\Pi$, $G_w=G_uG_v/\Pi$; and $y'=g'=-gy$, $Y'=G'=-GY$ in the gap variable.
\end{proposition}
\begin{proof}
For $i<j$ and $s=\sigma_j-\sigma_i$ we have
\[y=\frac{(1-d)Q_i}{Q_j-Q_i},\quad g=\frac{(1-d)Q_j}{Q_j-Q_i},\quad G=\frac{(1+d)P_i}{P_j-P_i},\quad Y=\frac{(1+d)P_j}{P_j-P_i}.\]
So $p=2\kappa_{ij}=d(1-K)$ and, as $y-g=d-1$ and $Y-G=1+d$, $m=Y-g=2d-p=d(1+K)$. The last sum of \eqref{eq:H} becomes the last bracket of \eqref{eq:gy}; for the middle point,
$(K(-u)-1)(K(v)-1)=(1+K(u))(1-K(v))$. Next,
\begin{align*}
Y_uy_v-g_uG_v&=(1-d^2)\lambda_b\Big(\frac1{(P_b-P_a)(Q_c-Q_b)}-\frac1{(Q_b-Q_a)(P_c-P_b)}\Big)=-\frac{(1-d^2)\lambda_bA}{D_{ab}D_{bc}},\\
\Pi&=\frac{(1+d)P_b(P_c-P_a)}{(P_b-P_a)(P_c-P_b)},\qquad\Theta=\frac{(1-d)Q_b(Q_c-Q_a)}{(Q_b-Q_a)(Q_c-Q_b)},
\end{align*}
using $A=(P_b-P_a)(Q_c-Q_b)-(P_c-P_b)(Q_b-Q_a)$. Hence $(Y_uy_v-g_uG_v)^2/(\Pi\Theta)=(1-d^2)\lambda_bA^2/(D_{ab}D_{bc}D_{ac})$, and since the first term of \eqref{eq:kappaform} is
$\frac{1-d^2}8A^2/(D_{ab}D_{bc}D_{ca})$, it contributes $12X_b(\dots)^2/(\Pi\Theta)$ to $96\hs$. The addition rules are immediate from the displayed fractions, for example
$y_uy_v/\Theta=(1-d)Q_a/(Q_c-Q_a)=y_w$. The derivative rules follow from $\partial_sB(k,s)=-B(B+k)$ with $y+(1-d)=g$ and $G+(1+d)=Y$.
\end{proof}

\section{The qubit inequality (62)}

\begin{theorem}\label{thm:62}
For every $\beta\in\R\setminus\{0,\frac12,1\}$, \eqref{62} holds strictly for $\hs_\beta$. In particular, for $n=2$ the scalar curvature $S_\beta$ increases strictly under mixing.
\end{theorem}
\begin{proof}
Let $y/x=e^{2s}$, $s>0$, $x+y=1$, so $\frac1x=2\cosh s\,e^{s}$ and $\frac1y=2\cosh s\,e^{-s}$. By \eqref{eq:twopoint} and $1-2\beta=-d=-\kappa_{xy}-\kappa_{yx}$,
\[\mathcal P=3\hs(x,x,y)+3\hs(x,y,y)=\frac12\Big(\frac{\kappa_{xy}^2-2d\kappa_{xy}}x+\frac{\kappa_{yx}^2-2d\kappa_{yx}}y\Big)=\frac12\Big(\frac{\kappa_{yx}^2-d^2}x+\frac{\kappa_{xy}^2-d^2}y\Big),\]
and with \eqref{eq:kappaK}, $\kappa_{xy}=\frac d2(1-K)$, $\kappa_{yx}=\frac d2(1+K)$, $K=K(s)$,
\begin{equation}\label{eq:S2}
S_\beta=\frac32+\frac{d^2}2\big(J^2-4\cosh^2s+1\big),\qquad J=K\cosh s+\sinh s .
\end{equation}
We show that $\mathcal G=J^2-4\cosh^2s$ has $\mathcal G'<0$ on $(0,\infty)$; as $s$ decreases when $x$ increases to $\frac12$, this is the condition in Proposition~\ref{prop:translate}(1). $\mathcal G$ depends on $d$ only
through $d^2$, so let $d>0$.

Let $W=\ph'\sinh s-\ph\cosh s$ and $\omega=W/\ph'=\sinh s-K\cosh s$. Then $W(0)=0$ and, by $\ph''=d^2\ph+\sinh s$, $W'=(\ph''-\ph)\sinh s=\frac1d\sinh(ds)\sinh s>0$. So $\omega>0$. Also $J=2\sinh s-\omega$, and
$J>0$ because $\ph,\ph'>0$; hence $0<\omega<2\sinh s$. A direct computation gives
\[\mathcal G'=-4\omega\cosh s-2\omega'(2\sinh s-\omega).\]
It remains to see that $\omega'\ge0$. $W$ and $\ph'$ vanish at $0$ and $\ph''>0$, so by the monotone form of l'H\^opital's rule $\omega$ is nondecreasing if $W'/\ph''$ is. For $d\ne1$,
\[\frac{W'}{\ph''}=\frac{1-d^2}{d\,q(s)},\qquad q(s)=\frac1{\sinh ds}-\frac d{\sinh s},\qquad q'(s)=d\,\big(\mathsf m(s)-\mathsf m(ds)\big),\quad \mathsf m(x)=\frac{\cosh x}{\sinh^2x},\]
with $\mathsf m$ strictly decreasing. If $d<1$ then $q>0$ and $q'<0$; if $d>1$ then $q<0$ and $q'>0$. In both cases $(1-d^2)/q$ is positive and increasing. (Signs of $q$: $\sinh s>d\sinh ds$ for $d<1$, and
the reverse for $d>1$.) Therefore $\mathcal G'\le-4\omega\cosh s<0$.
\end{proof}

\section{The inequality (63)}

\begin{theorem}\label{thm:63}
Let $0<d\le3$. Then $2\psi'+\psi''<0$ on $\R\setminus\{0\}$ and $e^{2s}\psi'$ is strictly decreasing on $\R$. In particular \eqref{63} holds for $\hs_\beta$ when $\beta\in[-1,2]\setminus\{0,\frac12,1\}$.
\end{theorem}
\noindent The one-variable statement is machine-checked in Lean (\texttt{lean63/Ineq63.lean}: \texttt{wyd\_ineq63}, and \texttt{bkm\_ineq63} for $d=1$). Its translation to \eqref{63} is Proposition~\ref{prop:translate}(2), proved by hand.
\begin{proof}
Let $s>0$ and $k=\ph/\ph'$, $S=\ph''/\ph'$, $v=1-kS$, $u=-v'$, $D=d^2$, $Y=(1-k^2)(1-Dk^2)$.

\emph{An autonomous system.} Let $J=(\ph'')^2-D(\ph')^2+(\ph')^2-D\ph^2$. Since $\ph''-D\ph=\sinh$ and $\ph'''-D\ph'=\cosh$, we have $J'=2\ph''\cosh s+2\ph'\sinh s=(2\ph'\cosh s)'$, so
$J=2\ph'\cosh s$. Dividing by $(\ph')^2$,
\begin{gather}
k'=v=1-kS,\qquad S'=\tfrac12(1+D-S^2-Dk^2),\label{eq:flow}\\
(S\mp1)^2-D(1\mp k)^2=2e^{\mp s}/\ph',\qquad\text{so}\quad e^{-2s}=\frac{(S-1)^2-D(1-k)^2}{(S+1)^2-D(1+k)^2}.\label{eq:E}
\end{gather}
From \eqref{eq:flow},
\begin{equation}\label{eq:u}
2ku=4v-3v^2-Y=(1-v)(3v-1)+k^2(1+D-Dk^2),\qquad u'=-2Su+v(1+D-2Dk^2).
\end{equation}

\emph{Signs.} $\ph,\ph',\ph''>0$, and $v>0$ because $(v\ph'^2)'=W>0$, with $W$ as in the proof of Theorem~\ref{thm:62}. As $s\to0^+$, $k\to0$, $v\to\frac13$ and $S\to\infty$. By \eqref{eq:E}, $S\ne1$, so $S>1$; then
$(\ph'-\ph)'=\ph''-\ph'>0$ gives $k<1$, and \eqref{eq:E} gives $S-1>d(1-k)$, i.e.
\begin{equation}\label{eq:vup}
0<v<(1-k)(1-dk),\qquad\text{in particular } dk<1.
\end{equation}
Since $2Dk^2\le1+D$, $(u\ph'^2)'=\ph'^2v(1+D-2Dk^2)\ge0$, and $u\ph'^2\to0$ at $0$; so $u\ge0$, and hence $v\le\frac13$.

\emph{Barrier principle.} If $f$ is differentiable on $(0,\infty)$, positive near $0$, and $f'>0$ wherever $f=0$, then $f\ge0$.

(Q) $v\ge\frac14(1-9k^2)$. Here $f=v-\frac14(1-9k^2)$ tends to $\frac1{12}$. If $f=0$ then $1-9k^2=4v>0$, and by \eqref{eq:u} and $D\le9$, $2kf'=2k(\frac92kv-u)\ge\frac1{16}(1-9k^2)(3-7k^2)>0$.

(G) $Y\le2(1+k)v+6v^2$. Here $f=2(1+k)v+6v^2-Y$ tends to $\frac13$. If $f=0$, eliminating $D$ and $u$ by \eqref{eq:u},
\[k(1-k^2)f'=v^2\big((42-30k^2)v-(1-3k)(7-5k)(1+k)\big).\]
This is positive for $k\ge\frac13$. For $k<\frac13$, (Q) bounds the bracket below by $\frac12(1-3k)(7+59k-5k^2-45k^3)>0$.

(C) $u\le2v$ if $k>\frac13$. By \eqref{eq:u}, $2k(2v-u)=3v^2-4(1-k)v+Y=:c(v)$. Let $\delta=dk$, so $Y=(1-k^2)(1-\delta^2)$. If $\delta\ge\frac13$, then $v<v^*=(1-k)(1-\delta)\le\frac23(1-k)$, where $c$ decreases, and
$c(v^*)=2(1-k)(1-\delta)k(1-d+2\delta)\ge0$, because $1-d+2dk\ge1-\frac d3\ge0$. If $\delta<\frac13$, then $c\ge(1-k)\big((1+k)(1-\delta^2)-\frac43(1-k)\big)\ge0$.

(M) $(1-k)v\ge e^{-2s}(3k-1)(1+k)$ if $k>\frac13$. By \eqref{eq:E} and $S=(1-v)/k$ this is
\[F:=(1-k)v\big((1+k-v)^2-Dk^2(1+k)^2\big)-(3k-1)(1+k)\big((1-k-v)^2-Dk^2(1-k)^2\big)\ge0 .\]
With $A=(1+k)v-(3k-1)(1-k)$,
\begin{align*}
F&=A\,(3v^2-4(1-k)v+Y)+v\big(2(3k-1)^2(1-k)+2k(5-7k)v-(2+4k)v^2\big)\\
&=-A\,(2(1+k)v+6v^2-Y)+v^2\big(7-22k+19k^2+(7+5k)v\big).
\end{align*}
If $A<0$ use the second form, (G), and $22^2<4\cdot7\cdot19$. If $A\ge0$ use the first form and (C): its last bracket $b(v)$ is concave in $v$ and $b(0)\ge0$; for $k\ge\frac12$, $v\le1-k$ and $b(1-k)=4k(1-k)(2k-1)\ge0$; for $k\le\frac12$, $v\le\frac13$ and
$b(\frac13)=-18k^3+\frac{76}3k^2-\frac{100}9k+\frac{16}9>0$ on $[\frac13,\frac12]$.

\emph{Negative side.} $K(-s)=-k$, $K'(-s)=v$, $K''(-s)=u$, and with $E=e^{-2s}$,
\begin{multline*}
-(2\psi'+\psi'')(-s)=8(1-k)v+4v^2+4(1-k)u\\+E\big(8(1-3k)(1+k)+12(3k+1)v+2(3k+1)u-6v^2\big).
\end{multline*}
Here $12(3k+1)v\ge6v^2$ and $u\ge0$, and $8(1-k)v+8E(1-3k)(1+k)\ge0$, by (M) if $k>\frac13$ and trivially otherwise; $4v^2>0$ is left.

\emph{Positive side.} $-E(2\psi'+\psi'')(s)=R_0+ER_1$ with $R_0=8(3k+1)(1-k)+12(1-3k)v-6v^2-2(1-3k)u$ and $R_1=8(k+1)v+4v^2-4(k+1)u$. If $k>\frac13$, then $R_1\ge0$ by (C), and $u\ge0$, $v\le1-k$ give $R_0\ge(1-k)(14-6k)>0$.
If $k\le\frac13$, \eqref{eq:u} with $v\le\frac13$ and $D\le9$ gives $u\le\frac12k(10-9k^2)$; then $R_0>0$ and $R_0+R_1>0$, because
$8(3k+1)(1-k)-\frac23-k(1-3k)(10-9k^2)$ and $8(3k+1)(1-k)-k(3-k)(10-9k^2)$ are positive on $[0,\frac13]$; and $0<E<1$.

Finally $K$ is $C^1$ at $0$, so $e^{2s}\psi'$ is continuous, with negative derivative off $0$.
\end{proof}

\section{The inequalities (64) and (65)}\label{sec:6465}

\begin{theorem}\label{thm:6465}
Let $0<d\le3$, $d\ne1$. Then, with $H$ as in Proposition~\ref{prop:H},
\begin{itemize}
\item[(b)] $(e^{-2\sigma_1}\partial_{\sigma_1}-e^{-2\sigma_2}\partial_{\sigma_2})H\ge0$ for distinct $\sigma_i$ with $\sigma_1<\sigma_2$,
\item[(c)] $(\partial_{\sigma_1}^2-2\partial_{\sigma_1})H\le0$ for distinct $\sigma_i$.
\end{itemize}
\end{theorem}
\begin{proof}
\emph{Newton expansion.} For real $E_1,\dots,E_k$ the divided difference of $E\mapsto e^{Eu}$,\[[E_1,\dots,E_k]e^{\,\cdot\,u}=u^{k-1}\int_{\Delta_{k-1}}e^{u\sum t_iE_i}\,dt,\]is positive for $u>0$, in whatever order, coincident nodes included. If
$f(u)=\sum_Ec_Ee^{Eu}$ with all $E$ among $E_1,\dots,E_k$, then
\begin{equation}\label{eq:newton}
f(u)=\sum_{i=1}^kf_{i}(0)\,[E_1,\dots,E_i]e^{\,\cdot\,u},\qquad f_1=f,\quad f_{i+1}=f_i'-E_if_i,\qquad f_i(0)=\sum_Ec_E\prod_{l<i}(E-E_l),
\end{equation}
so $f\ge0$ on $[0,\infty)$ if every $f_i(0)\ge0$. (Directly: $(e^{-E_iu}f_i)'=e^{-E_iu}f_{i+1}$ and $f_{k+1}=0$.) For $f(u,v)=\sum c\,e^{Eu+Fv}$ the same holds on the quadrant with the coefficients $f_{ij}(0,0)=\sum c\prod_{l<i}(E-E_l)\prod_{l<j}(F-F_l)$.

\emph{The numerators.} Let $p,q>0$ with $d=q/p$, $F(s)=p\sinh ds-q\sinh s$ and $G(s)=q(\cosh ds-\cosh s)$. By Proposition~\ref{prop:H},
\begin{gather*}
H=\frac{A}{24p^2\prod2G},\qquad A=3q(q^2-p^2)N^2-2p^2\sum_{\rm cyc}e^{-2\sigma_1}\Phi(\sigma_2-\sigma_1)\Phi(\sigma_3-\sigma_1)\,2G(\sigma_2-\sigma_3),
\end{gather*}
$\Phi=2F-2G$, $N=\sum_{\rm cyc}e^{-\sigma_3}2F(\sigma_1-\sigma_2)$, $\prod 2G$ over the three pairs. Order the points, with gaps $u,v>0$. Clearing denominators in (b) and (c), for each of the three positions of the third (resp.\ differentiated) variable, leaves
a positive exponential times $\sum c\,e^{(i+jd)u+(i'+j'd)v}$ with $c\in\mathbb Z[p,q]$ homogeneous; the cleared denominator is a square in (b), and in (c) it has the sign of $2G$, which is that of $d-1$. The six numerators have the following sizes.
\begin{center}
\begin{tabular}{lcccccc}
\toprule
& (c), 1 & (c), 2 & (c), 3 & (b), 1 & (b), 2 & (b), 3\\
\midrule
terms & 616 & 616 & 616 & 554 & 554 & 554\\
exponents in $u$, $v$ & 55, 28 & 28, 28 & 28, 55 & 28, 36 & 36, 36 & 36, 28\\
\bottomrule
\end{tabular}
\end{center}

\emph{The check.} For $0<d<1$ put $p=1+t$, $q=t$, $t=d/(1-d)>0$. For $1<d\le3$ put $p=1+t$, $q=3+t$, $t=(3-d)/(d-1)\ge0$. Replacing $(u,v)$ by $(pu,pv)$, the exponents $ip+jq$ and the coefficients $c$ are polynomials in $t$ with integer coefficients, so
the Newton coefficients are too. Use one fixed order of the nodes in each range: increasing at $d=0.501$, resp.\ at $d=2.001$. In all twelve cases, \emph{every Newton coefficient is a polynomial in $t$ with
nonnegative integer coefficients} ($10\,928$ nonzero Newton coefficients, $470\,258$ integer coefficients counted from the lowest power, no negative one). Hence the numerators are nonnegative.
\end{proof}

The check is integer arithmetic and takes about two seconds in Python. It is sharp: just above $d=3$ negative Newton coefficients appear in (c) 2, (c) 3 and (b) 1, the cases that degenerate at $d=3$.

\section{Conclusion and verification}\label{sec:status}

Combining Theorems~\ref{thm:closed}, \ref{thm:reduction}, \ref{thm:62}, \ref{thm:63}, \ref{thm:6465} and Proposition~\ref{prop:translate}:

\begin{corollary}\label{cor:main}
Let $\beta\in[-1,2]\setminus\{0,\frac12,1\}$ and $n\ge2$. The scalar curvature of $(\Dn,g^\beta)$ increases strictly under mixing: $S_\beta(\varrho)<S_\beta(\rho)$ if $\rho$ is more mixed than $\varrho$ and the spectra differ.
\end{corollary}

For $\beta=\frac12$, $\hs_\beta=0$ and $S_\beta$ is constant. $S_\beta$ depends on $\beta$ only through $d^2$, so $d\in(0,3]$, $d\ne1$, covers $\beta\in[-1,2]\setminus\{0,\frac12,1\}$. The limits $\beta\to0,1$ are the Kubo--Mori case (Petz's conjecture).

\subsection*{Lean} The folder \texttt{lean/WYD} contains a Lean~4 proof of Corollary~\ref{cor:main} with no hypothesis beyond those of the statement (\texttt{wyd\_curvature\_lt\_of\_isMajorizedBy}). It includes
Theorem~\ref{thm:closed} for all $\beta\ne0,1$ (\texttt{scalarCurvature\_wydMetric}), the reduction for any symmetric smooth $\hs$, \eqref{62}--\eqref{65} for $\beta\in[-1,2]\setminus\{0,\frac12,1\}$
(\texttt{wyd\_62}, \dots, \texttt{wyd\_65}), Proposition~\ref{prop:H} (\texttt{wydHs\_eq\_hypH}) and the parts of Proposition~\ref{prop:translate} that are used. The metric is defined as the pullback of a constant pairing by
$\rho\mapsto(\rho^\beta,\rho^{1-\beta})$, which is \eqref{eq:metric} (\texttt{wydMetric\_apply}, \texttt{wydMetric\_diag\_spectral}). Only the axioms \texttt{propext},
\texttt{Classical.choice} and \texttt{Quot.sound} are used, and there is no \texttt{native\_decide}: the twelve checks of Theorem~\ref{thm:6465} are evaluated by the kernel (\texttt{C64\_*}, \texttt{C65\_*}), including the construction of the numerators, and the
Newton argument is proved in \texttt{NewtonSound}, \texttt{NewtonEP}, \texttt{NewtonNum}. The Lean proof of \eqref{62} replaces l'H\^opital's rule by $\omega'=k(S\cosh s-\sinh s)>0$, with $k,S$ as in Section~7 and $S>1$. Not in Lean: Proposition~\ref{prop:gy}, which is no longer used.

\subsection*{The conjecture with the definitions of \texorpdfstring{\cite{gi}}{the paper}} The file \texttt{Literal.lean} states Conjecture~8.1 of \cite{gi} with the definitions of that paper and proves it (\texttt{conjecture81}): the density matrices (\texttt{density}, Definition~2.5),
the order by partial sums of the decreasing rearrangements (\texttt{MoreMixed}, Definitions~2.3 and~2.6), Schur-increasing functions (\texttt{SchurIncreasingOn}, Definition~2.8, non-strict, spectra that coincide up to order included), $f_p$ of Definition~6.1, and the metric
$\operatorname{Tr}(A\,c_{f_p}(L_\rho,R_\rho)(B))$ on the traceless Hermitian matrices (\texttt{IsWYDMetric}, Definition~6.2); $c_{f}(L_\rho,R_\rho)$ is the linear map that multiplies $uv^*$ by $c_f(a,b)$ when $\rho u=au$, $\rho v=bv$ (\texttt{cLR\_add}, \texttt{cLR\_smul}, \texttt{cLR\_vecMulVec}).
The scalar curvature at $\rho$ is computed in the affine chart at $\rho$; any other affine chart gives the same (\texttt{scal\_eq\_chart}). The scalar curvature is Schur-increasing for $p\le-1$ and for $p\ge\frac12$, $p\ne1$, and every $n$ (\texttt{wyd\_schurIncreasing}); so any $\varepsilon>0$ will do. It is strictly Schur-increasing if moreover $p\ne2$ (\texttt{wyd\_strictly}).
\emph{The case $p=1$} (the same metric as $p=\infty$; $f_1=(x-1)/\log x$, the Kubo--Mori metric) is Petz's conjecture, proved in the PetzCurvature library. Its modules are rebuilt in \texttt{lean/KM} on the same shared files as the rest (only imports change), and \texttt{KuboMori.lean} shows that its BKM form $\int_0^1\operatorname{Re}\operatorname{tr}(A R B R)\,ds$, $R=((1-s)\rho+s)^{-1}$, is the metric of Definition~6.2 at $p=1$ at every state (\texttt{bkmForm\_eq}). So the scalar curvature is Schur-increasing for every $p\le-1$ and every $p\ge\frac12$ (\texttt{wyd\_schurIncreasing\_all}), strictly for $p\ne2$ (\texttt{wyd\_strictly\_all}); at $p=1$ its value at $I/n$ is Petz's $\frac18(n^2-1)(n^2-4)$ (\texttt{scal\_kuboMori\_uniform}). That these $p$ are exactly the ones for which $f_p$ is operator monotone (Theorem~6.1 of \cite{gi}) is also formalized (\texttt{Theorem61.lean}, \texttt{theorem61}, with the definition of \cite{gi} over $n\times n$ matrices, normalized and symmetric): for $p\ge1$ and $p\le-1$, $1/f_p$ is an average over $[0,1]^2$ of inverses of operator geometric means of the operator monotone functions $1-s+sx$ and $x/(1-s+sx)$, so $f_p$ is operator monotone in every unital C$^*$-algebra; $\frac12\le p<1$ follows from $f_p=f_{\tilde p}$; for $p\in(0,\frac12)$, $f_p(x)\to0$ as $x\to\infty$, and $p\in(-1,0)$ reduces to it. Controls in the same file: the metric exists, is symmetric and positive definite at every state (\texttt{exists\_isWYDMetric}, \texttt{IsWYDMetric.symm}, \texttt{IsWYDMetric.pos});
at $p=2$ the curvature is $\frac14(n^2-1)(n^2-2)$ (\texttt{scal\_wigner\_yanase}), which fixes the normalization and shows that strictness fails there; for $p\ne1,2$ the state $I/n$ has strictly larger curvature than every other one (\texttt{scal\_lt\_uniform}), so the hypotheses can be met and the direction is the right one.

\subsection*{Earlier computer-assisted proof} Before Theorem~\ref{thm:6465}, (b) and (c) were certified by interval arithmetic for every $d\in[0,3]$ (\texttt{NOTES.md}, Steps 12--14; $351\,624$ Arb boxes, no failures; Proposition~\ref{prop:gy} is the form used for the tails). That computation is now
redundant and serves only as an independent check, as do Steps 9, 10 and 15 for \eqref{63}.

\subsection*{Review} None of this has been independently reviewed by a person. Lean checks the proof of the formal statement; that the formal statement is the conjecture still has to be read. The shortest way is \texttt{Literal.lean}, whose definitions are those of \cite{gi}; what it takes from elsewhere is the scalar curvature of a metric in a chart (\texttt{AlphaCurvature/Curvature/Basic.lean}) and the space of traceless Hermitian matrices.
At random non-diagonal states ($n=2,3$) the Lean definitions, the metric with $f_p$, Dittmann's formula (Theorem~5.3 of \cite{gi}) and \eqref{eq:closed} were computed separately and agree (\texttt{fidelity/fidelity.py}); a search for counterexamples by descent (1800 runs, $n\le6$) found none, while the same search finds one for the Bures metric in 36 of 36 runs (\texttt{fidelity/adversary.py}). All compiled files, and the AlphaCurvature files they use, were replayed through the kernel (\texttt{leanchecker}'s replay, one module at a time, on Lean 4.36.0: 37 WYD, 24 AlphaCurvature and 20 \texttt{KM} modules, no failures; not Mathlib). A false theorem compiled with the kernel switched off is rejected by the same replay.
The displayed identities of Sections 3, 5 and 6 have also been checked by a script in the \texttt{checks} folder, \texttt{reduction\_check.py} (27 checks: the single-gap and divided-difference identities symbolically in
$\beta$, the hyperbolic forms at 300 random points to $10^{-55}$ or better, the contraction formula on a random orthonormal basis, and \eqref{eq:closed} against brute-force curvature). The numerators of Theorem~\ref{thm:6465} agree with an independent evaluator at 120 random points to $10^{-135}$, and an independent implementation of the check finds the same signs and detects 36 of 36 planted coefficient errors.

\begin{thebibliography}{9}
\bibitem{bhatia} R.~Bhatia, \emph{Matrix Analysis}, Springer, 1997. (Chapters II and V: majorization; Daleckii--Krein formula.)
\bibitem{dittmann} J.~Dittmann, \emph{On the curvature of monotone metrics and a conjecture concerning the Kubo--Mori metric}, Linear Algebra Appl.\ 315 (2000) 83--112, arXiv:quant-ph/9906009.
\bibitem{gi} P.~Gibilisco and T.~Isola, \emph{On the monotonicity of scalar curvature in classical and quantum information geometry}, J.\ Math.\ Phys.\ 46 (2005), arXiv:math-ph/0407007. (Conjecture 8.1.)
\bibitem{hp} H.~Hasegawa and D.~Petz, \emph{Non-commutative extension of information geometry II}, in: Quantum Communication, Computing, and Measurement, Plenum, 1997. (Only for the range $-1\le\beta\le2$; not used in the proofs.)
\bibitem{shima} H.~Shima, \emph{The Geometry of Hessian Structures}, World Scientific, 2007.
\end{thebibliography}

\end{document}
